\chapter{Ejemplos Detallados de Intervenciones}
\label{anexo:ejemplos-intervenciones}

Este anexo presenta ejemplos detallados de intervenciones para diferentes categorías de \gls{community-smells}. Cada ejemplo incluye diagnóstico, estrategia paso a paso, métricas de éxito y resultados observados.

\section{Intervención: \textit{Radio Silence} (\textit{Smell} de Comunicación)}

\begingroup
\scriptsize
\renewcommand{\arraystretch}{1.4}
\setlength{\tabcolsep}{4pt}
\setlength{\LTleft}{0pt}
\setlength{\LTright}{0pt}

\begin{longtable}{
    >{\RaggedRight\hspace{0pt}}p{3.5cm}
    >{\RaggedRight\hspace{0pt}}p{12.5cm}
    }

    \caption{Ejemplo de intervención: \textit{Radio Silence}} \label{tab:radio-silence}                          \\
    \toprule
    \multicolumn{2}{l}{\textbf{Intervención: \textit{Radio Silence}}}                                            \\
    \midrule
    \endfirsthead

    \multicolumn{2}{c}{{\bfseries \tablename\ \thetable{} -- Continuación}}                                      \\
    \toprule
    \multicolumn{2}{l}{\textbf{Intervención: \textit{Radio Silence}}}                                            \\
    \midrule
    \endhead

    \midrule
    \multicolumn{2}{r}{\textit{Continúa en la siguiente página...}}                                              \\
    \endfoot

    \bottomrule
    \endlastfoot

    \textbf{Equipo}              & Alpha (7 personas)                                                            \\
    \textbf{Severidad detectada} & Moderada (afecta la velocidad, no bloquea)                                    \\
    \midrule

    \multicolumn{2}{l}{\textbf{Diagnóstico}}                                                                     \\* \nopagebreak
                                 & \textit{SNA} muestra que María centraliza el 70~\% de la comunicación         \\
                                 & 3 \textit{devs} reportan ``enterarse tarde'' de cambios en la \textit{API}    \\
                                 & En la \textit{Daily}, solo María y Pedro hablan regularmente                  \\
                                 & \textit{Health Score} ``Comunicación'' 4,5/10                                 \\
    \midrule

    \multicolumn{2}{l}{\textbf{Estrategia (5 SSP)}}                                                              \\* \nopagebreak
    \multicolumn{2}{l}{\textit{Semana 1:}}                                                                       \\* \nopagebreak
                                 & \textit{Workshop} ``Comunicación distribuida'' (2h)                           \\
                                 & Crear canal \#alpha-actualizaciones para cambios técnicos                     \\
                                 & Norma: quien cambia una \textit{API} pública, lo publica en el canal          \\
    \multicolumn{2}{l}{\textit{Semana 2:}}                                                                       \\* \nopagebreak
                                 & Rotación de ``facilitador de \textit{Daily}'' (cada día alguien diferente)    \\
                                 & Pregunta explícita en la \textit{Daily}: ``¿Hay algo que otros deban saber?'' \\
    \midrule

    \multicolumn{2}{l}{\textbf{Métricas de éxito}}                                                               \\* \nopagebreak
                                 & Mensajes en \#alpha-actualizaciones: >= 5/semana                              \\
                                 & Participación en la \textit{Daily}: todos hablan al menos 3x/semana           \\
                                 & Encuesta: ``Me entero a tiempo'' sube de 4/10 a 7/10                          \\
    \midrule

    \multicolumn{2}{l}{\textbf{Resultado (Sprint +1)}}                                                           \\* \nopagebreak
                                 & Mensajes en canal: 8/semana [OK]                                              \\
                                 & Participación \textit{Daily}: 6 de 7 hablan regularmente [OK]                 \\
                                 & Encuesta: 6,5/10 (mejora parcial, continuar)                                  \\
\end{longtable}
\endgroup

\section{Intervención: \textit{Truck Factor} = 1 (\textit{Smell} de Conocimiento)}

\begingroup
\scriptsize
\renewcommand{\arraystretch}{1.4}
\setlength{\tabcolsep}{4pt}
\setlength{\LTleft}{0pt}
\setlength{\LTright}{0pt}

\begin{longtable}{
    >{\RaggedRight\hspace{0pt}}p{3.5cm}
    >{\RaggedRight\hspace{0pt}}p{12.5cm}
    }

    \caption{Ejemplo de intervención: \textit{Truck Factor} = 1} \label{tab:truck-factor}                             \\
    \toprule
    \multicolumn{2}{l}{\textbf{Intervención: \textit{Truck Factor} = 1}}                                              \\
    \midrule
    \endfirsthead

    \multicolumn{2}{c}{{\bfseries \tablename\ \thetable{} -- Continuación}}                                           \\
    \toprule
    \multicolumn{2}{l}{\textbf{Intervención: \textit{Truck Factor} = 1}}                                              \\
    \midrule
    \endhead

    \midrule
    \multicolumn{2}{r}{\textit{Continúa en la siguiente página...}}                                                   \\
    \endfoot

    \bottomrule
    \endlastfoot

    \textbf{Equipo}             & Beta (5 personas)                                                                   \\
    \textbf{Componente crítico} & Motor de reglas de negocio                                                          \\
    \textbf{Único experto}      & Carlos                                                                              \\
    \midrule

    \multicolumn{2}{l}{\textbf{Diagnóstico}}                                                                          \\* \nopagebreak
                                & Carlos escribió el 95~\% del código del motor                                       \\
                                & Cuando Carlos está de vacaciones, los \textit{bugs} en el motor quedan sin resolver \\
                                & Otros \textit{devs} ``tienen miedo'' de tocar ese código                            \\
    \midrule

    \multicolumn{2}{l}{\textbf{Estrategia (8 SSP distribuidos en 3 \textit{Sprints})}}                                \\* \nopagebreak
    \multicolumn{2}{l}{\textit{Sprint 1 (3 SSP):}}                                                                    \\* \nopagebreak
                                & Carlos documenta la arquitectura en 3 páginas                                       \\
                                & Carlos hace \gls{pair-programming} con Ana en 2 tareas del motor                    \\
    \multicolumn{2}{l}{\textit{Sprint 2 (3 SSP):}}                                                                    \\* \nopagebreak
                                & Ana hace tarea sola en el motor, Carlos solo revisa                                 \\
                                & Carlos hace \gls{pair-programming} con Pedro en 1 tarea                             \\
                                & Sesión de 2h: ``Cómo depurar (\textit{debuggear}) el motor''                        \\
    \multicolumn{2}{l}{\textit{Sprint 3 (2 SSP):}}                                                                    \\* \nopagebreak
                                & Pedro hace tarea solo, Carlos revisa                                                \\
                                & Documentar la guía de resolución de problemas (\textit{troubleshooting guide})      \\
                                & Evaluar: ¿Ana y Pedro pueden resolver \textit{bugs} nivel 1-2?                      \\
    \midrule

    \multicolumn{2}{l}{\textbf{Métricas de éxito}}                                                                    \\* \nopagebreak
                                & \textit{Truck factor}: de 1 a 3 (Carlos + Ana + Pedro)                              \\
                                & Tiempo de resolución de \textit{bugs} cuando Carlos no está: < 48h                  \\
                                & Ana y Pedro reportan confianza >= 6/10 para trabajar en el motor                    \\
    \midrule

    \multicolumn{2}{l}{\textbf{Resultado}}                                                                            \\* \nopagebreak
                                & \textit{Sprint} 3: Ana resolvió un \textit{bug} de producción sola [OK]             \\
                                & Confianza: Ana 7/10, Pedro 5/10 (continuar con Pedro)                               \\
                                & \textit{Truck factor} efectivo: 2,5 (Ana casi independiente, Pedro parcial)         \\
\end{longtable}
\endgroup

\section{Intervención: Conflicto No Resuelto (\textit{Smell} Relacional)}

\begingroup
\scriptsize
\renewcommand{\arraystretch}{1.4}
\setlength{\tabcolsep}{4pt}
\setlength{\LTleft}{0pt}
\setlength{\LTright}{0pt}

\begin{longtable}{
    >{\RaggedRight\hspace{0pt}}p{3.5cm}
    >{\RaggedRight\hspace{0pt}}p{12.5cm}
    }

    \caption{Ejemplo de intervención: Conflicto No Resuelto} \label{tab:conflicto-no-resuelto}                                       \\
    \toprule
    \multicolumn{2}{l}{\textbf{Intervención: Conflicto No Resuelto}}                                                                 \\
    \midrule
    \endfirsthead

    \multicolumn{2}{c}{{\bfseries \tablename\ \thetable{} -- Continuación}}                                                          \\
    \toprule
    \multicolumn{2}{l}{\textbf{Intervención: Conflicto No Resuelto}}                                                                 \\
    \midrule
    \endhead

    \midrule
    \multicolumn{2}{r}{\textit{Continúa en la siguiente página...}}                                                                  \\
    \endfoot

    \bottomrule
    \endlastfoot

    \textbf{Equipo} & Gamma (6 personas)                                                                                             \\
    \textbf{Partes} & Dev A (\textit{frontend}) vs Dev B (\textit{backend})                                                          \\
    \textbf{Origen} & Desacuerdo sobre diseño de \textit{API} hace 2 meses                                                           \\
    \midrule

    \multicolumn{2}{l}{\textbf{Diagnóstico}}                                                                                         \\* \nopagebreak
                    & Las \textit{code reviews} entre A y B son hostiles                                                             \\
                    & Otros miembros evitan tareas que requieren A+B                                                                 \\
                    & La \textit{Daily} tiene tensión cuando ambos reportan                                                          \\
    \midrule

    \multicolumn{2}{l}{\textbf{Estrategia (5 SSP)}}                                                                                  \\* \nopagebreak
    \multicolumn{2}{l}{\textit{Día 1 -- Sesión 1:1 con Dev A (1h):}}                                                                 \\* \nopagebreak
                    & Escuchar su perspectiva sin juzgar                                                                             \\
                    & Identificar: ¿qué necesita para mejorar la situación?                                                          \\
                    & A dice: ``B nunca acepta mis sugerencias, siempre impone''                                                     \\
    \multicolumn{2}{l}{\textit{Día 2 -- Sesión 1:1 con Dev B (1h):}}                                                                 \\* \nopagebreak
                    & Escuchar su perspectiva sin juzgar                                                                             \\
                    & B dice: ``A no entiende las restricciones del \textit{backend}''                                               \\
    \multicolumn{2}{l}{\textit{Día 3 -- Preparación:}}                                                                               \\* \nopagebreak
                    & Identificar intereses comunes (ambos quieren un producto exitoso)                                              \\
                    & Preparar la estructura para la sesión conjunta                                                                 \\
    \multicolumn{2}{l}{\textit{Día 4 -- Sesión conjunta facilitada (2h):}}                                                           \\* \nopagebreak
                    & Reglas: hablar en primera persona, no interrumpir                                                              \\
                    & Cada uno expresa su perspectiva (10 min c/u)                                                                   \\
                    & Identificar puntos de acuerdo                                                                                  \\
                    & Negociar: ¿cómo manejar desacuerdos futuros?                                                                   \\
                    & Documentar acuerdos firmados por ambos                                                                         \\
    \multicolumn{2}{l}{\textit{Día 15 -- Seguimiento (30 min):}}                                                                     \\* \nopagebreak
                    & ¿Se están respetando los acuerdos?                                                                             \\
                    & Ajustar si es necesario                                                                                        \\
    \midrule

    \multicolumn{2}{l}{\textbf{Acuerdos documentados}}                                                                               \\* \nopagebreak
                    & 1. Desacuerdos técnicos se discuten en 1:1, no en comentarios del \textit{Pull Request} (\textit{PR comments}) \\
                    & 2. Si no hay acuerdo en 30 min, escalar al \textit{Tech Lead}                                                  \\
                    & 3. Ambos se comprometen a asumir buena intención del otro                                                      \\
    \midrule

    \multicolumn{2}{l}{\textbf{Métricas de éxito}}                                                                                   \\* \nopagebreak
                    & \textit{Code reviews} sin comentarios hostiles: 4 semanas consecutivas                                         \\
                    & A y B aceptan trabajar en tarea conjunta sin resistencia                                                       \\
                    & \textit{Health Score} ``Respeto'' sube de 4/10 a 6/10                                                          \\
    \midrule

    \multicolumn{2}{l}{\textbf{Resultado (Sprint +2)}}                                                                               \\* \nopagebreak
                    & 0 comentarios hostiles en los PR [OK]                                                                          \\
                    & Aceptaron la tarea conjunta en el \textit{Sprint} 16 [OK]                                                      \\
                    & Respeto: 6,5/10 [OK]                                                                                           \\
\end{longtable}
\endgroup

\section{Intervención: \textit{Meeting Hell} (\textit{Smell} Organizacional)}

\begingroup
\scriptsize
\renewcommand{\arraystretch}{1.4}
\setlength{\tabcolsep}{4pt}
\setlength{\LTleft}{0pt}
\setlength{\LTright}{0pt}

\begin{longtable}{
    >{\RaggedRight\hspace{0pt}}p{3.5cm}
    >{\RaggedRight\hspace{0pt}}p{12.5cm}
    }

    \caption{Ejemplo de intervención: \textit{Meeting Hell}} \label{tab:meeting-hell}                                 \\
    \toprule
    \multicolumn{2}{l}{\textbf{Intervención: \textit{Meeting Hell}}}                                                  \\
    \midrule
    \endfirsthead

    \multicolumn{2}{c}{{\bfseries \tablename\ \thetable{} -- Continuación}}                                           \\
    \toprule
    \multicolumn{2}{l}{\textbf{Intervención: \textit{Meeting Hell}}}                                                  \\
    \midrule
    \endhead

    \midrule
    \multicolumn{2}{r}{\textit{Continúa en la siguiente página...}}                                                   \\
    \endfoot

    \bottomrule
    \endlastfoot

    \textbf{Equipo}   & Delta (8 personas)                                                                            \\
    \textbf{Problema} & Promedio de 25h/semana en reuniones por persona                                               \\
    \midrule

    \multicolumn{2}{l}{\textbf{Diagnóstico}}                                                                          \\* \nopagebreak
    \multicolumn{2}{l}{\textit{Auditoría de calendarios muestra:}}                                                    \\* \nopagebreak
                      & 8h/semana: reuniones de Scrum (normal)                                                        \\
                      & 10h/semana: reuniones con partes interesadas (\textit{stakeholders}) (excesivo)               \\
                      & 7h/semana: reuniones ``informativas'' donde el equipo es pasivo                               \\
                      & Foco del equipo: 4,2/10 (crítico)                                                             \\
                      & Velocidad cayó 30~\% en los últimos 3 \textit{Sprints}                                        \\
    \midrule

    \multicolumn{2}{l}{\textbf{Estrategia (5 SSP)}}                                                                   \\* \nopagebreak
    \multicolumn{2}{l}{\textit{Semana 1 -- Documentación:}}                                                           \\* \nopagebreak
                      & Crear una hoja de cálculo (\textit{spreadsheet}): reunión, convocador, frecuencia, asistentes \\
                      & Calcular costo: horas $\times$ personas $\times$ tarifa hora $\approx$ \$X/semana             \\
    \multicolumn{2}{l}{\textit{Semana 2 -- Propuesta:}}                                                               \\* \nopagebreak
                      & Clasificar reuniones: eliminar, reducir frecuencia, o delegar                                 \\
                      & Propuesta: representante único para reuniones informativas                                    \\
                      & Establecer ``\textit{Focus Fridays}'': sin reuniones externas                                 \\
    \multicolumn{2}{l}{\textit{Semana 3 -- Negociación:}}                                                             \\* \nopagebreak
                      & Reunión con \textit{stakeholders} clave (SM lidera)                                           \\
                      & Presentar datos del impacto en la velocidad                                                   \\
                      & Acordar: representante rota, resumen escrito después                                          \\
    \multicolumn{2}{l}{\textit{Semana 4 -- Implementación:}}                                                          \\* \nopagebreak
                      & \textit{Focus Fridays} en efecto                                                              \\
                      & Representante único para 5 reuniones                                                          \\
                      & Monitorear calendarios                                                                        \\
    \midrule

    \multicolumn{2}{l}{\textbf{Métricas de éxito}}                                                                    \\* \nopagebreak
                      & Horas en reuniones: de 25h a 15h/semana                                                       \\
                      & Foco: de 4,2 a 6,0                                                                            \\
                      & Velocidad: recuperar al menos el 50~\% de la caída                                            \\
    \midrule

    \multicolumn{2}{l}{\textbf{Resultado (Sprint +2)}}                                                                \\* \nopagebreak
                      & Reuniones: 16h/semana (36~\% de reducción) [OK]                                               \\
                      & Foco: 5,8/10 (mejora significativa) [OK]                                                      \\
                      & Velocidad: +20~\% (recuperación parcial, continuar)                                           \\
\end{longtable}
\endgroup

\section{Experimento: Reducir \textit{Knowledge Islands}}

\begingroup
\scriptsize
\renewcommand{\arraystretch}{1.4}
\setlength{\tabcolsep}{4pt}
\setlength{\LTleft}{0pt}
\setlength{\LTright}{0pt}

\begin{longtable}{
    >{\RaggedRight\hspace{0pt}}p{3.5cm}
    >{\RaggedRight\hspace{0pt}}p{12.5cm}
    }

    \caption{Experimento: Reducir \textit{Knowledge Islands}} \label{tab:knowledge-islands}                                                                                                                                                                                                                                                  \\
    \toprule
    \multicolumn{2}{l}{\textbf{Experimento: Reducir \textit{Knowledge Islands}}}                                                                                                                                                                                                                                                             \\
    \midrule
    \endfirsthead

    \multicolumn{2}{c}{{\bfseries \tablename\ \thetable{} -- Continuación}}                                                                                                                                                                                                                                                                  \\
    \toprule
    \multicolumn{2}{l}{\textbf{Experimento: Reducir \textit{Knowledge Islands}}}                                                                                                                                                                                                                                                             \\
    \midrule
    \endhead

    \midrule
    \multicolumn{2}{r}{\textit{Continúa en la siguiente página...}}                                                                                                                                                                                                                                                                          \\
    \endfoot

    \bottomrule
    \endlastfoot

    \multicolumn{2}{l}{\textbf{Contexto}}                                                                                                                                                                                                                                                                                                    \\* \nopagebreak
                                      & Equipo Zeta tiene 3 subgrupos: \textit{Frontend}, \textit{Backend}, \textit{Data}.                                                                                                                                                                                                                   \\
                                      & Casi no hay comunicación entre subgrupos.                                                                                                                                                                                                                                                            \\
                                      & La integración al final del \textit{Sprint} genera conflictos.                                                                                                                                                                                                                                       \\
                                      & \textit{Health Score} ``Compromiso'': 5,5/10                                                                                                                                                                                                                                                         \\
    \midrule

    \multicolumn{2}{l}{\textbf{Hipótesis}}                                                                                                                                                                                                                                                                                                   \\* \nopagebreak
                                      & ``Si implementamos \textit{code reviews} cruzados obligatorios (1 PR/semana revisado por miembro de otro subgrupo), entonces el conocimiento compartido aumentará y los conflictos de integración disminuirán, porque la exposición temprana al código de otros reduce sorpresas y genera empatía.'' \\
    \midrule

    \multicolumn{2}{l}{\textbf{Métricas}}                                                                                                                                                                                                                                                                                                    \\* \nopagebreak
    \textit{Métrica}                  & \textit{Línea base $\rightarrow$ Target}                                                                                                                                                                                                                                                             \\* \nopagebreak
    PRs con \textit{reviewer} cruzado & 0/semana $\rightarrow$ 5/semana                                                                                                                                                                                                                                                                      \\
    Conflictos de integración         & 4/\textit{Sprint} $\rightarrow$ 2/\textit{Sprint}                                                                                                                                                                                                                                                    \\
    ``Entiendo el código de otros''   & 3/10 $\rightarrow$ 5/10                                                                                                                                                                                                                                                                              \\
    \textit{Health Score} Compromiso  & 5,5 $\rightarrow$ 6,5                                                                                                                                                                                                                                                                                \\
    \midrule

    \multicolumn{2}{l}{\textbf{Intervención}}                                                                                                                                                                                                                                                                                                \\* \nopagebreak
    \multicolumn{2}{l}{\textit{Semana 1:}}                                                                                                                                                                                                                                                                                                   \\* \nopagebreak
                                      & \textit{Workshop}: ``Por qué hacer \textit{code reviews} cruzados'' (1h)                                                                                                                                                                                                                             \\
                                      & Configurar GitHub: \textit{reviewer} aleatorio de otro subgrupo                                                                                                                                                                                                                                      \\
                                      & Norma: al menos 1 PR cruzado por persona/semana                                                                                                                                                                                                                                                      \\
    \multicolumn{2}{l}{\textit{Semanas 2--4:}}                                                                                                                                                                                                                                                                                               \\* \nopagebreak
                                      & Monitorear cumplimiento                                                                                                                                                                                                                                                                              \\
                                      & 1:1 con personas que no participan                                                                                                                                                                                                                                                                   \\
                                      & Celebrar en la \textit{Daily} cuando una revisión cruzada genera \textit{insight}                                                                                                                                                                                                                    \\
    \midrule

    \textbf{Duración}                 & 2 \textit{Sprints}                                                                                                                                                                                                                                                                                   \\
    \midrule

    \multicolumn{2}{l}{\textbf{Criterios de evaluación}}                                                                                                                                                                                                                                                                                     \\* \nopagebreak
                                      & Éxito: Conflictos <= 2, ``Entiendo código'' >= 5                                                                                                                                                                                                                                                     \\
                                      & Parcial: Conflictos <= 3, alguna mejora en la percepción                                                                                                                                                                                                                                             \\
                                      & Fracaso: Sin mejora o resistencia activa del equipo                                                                                                                                                                                                                                                  \\
    \midrule

    \multicolumn{2}{l}{\textbf{Resultados (Sprint +2)}}                                                                                                                                                                                                                                                                                      \\* \nopagebreak
                                      & PRs cruzados: 6/semana [OK] (superó el \textit{target})                                                                                                                                                                                                                                              \\
                                      & Conflictos: 3/\textit{Sprint} (mejora, no alcanzó el \textit{target})                                                                                                                                                                                                                                \\
                                      & ``Entiendo el código'': 4,5/10 (mejora parcial)                                                                                                                                                                                                                                                      \\
                                      & Compromiso: 6,0/10 (mejora)                                                                                                                                                                                                                                                                          \\
    \midrule

    \textbf{Evaluación}               & Éxito parcial                                                                                                                                                                                                                                                                                        \\
    \midrule

    \multicolumn{2}{l}{\textbf{Aprendizajes}}                                                                                                                                                                                                                                                                                                \\* \nopagebreak
                                      & 1. Las revisiones cruzadas toman más tiempo (factor 1,5x)                                                                                                                                                                                                                                            \\
                                      & 2. El sentido \textit{Backend} $\rightarrow$ \textit{Frontend} funciona mejor que el inverso                                                                                                                                                                                                         \\
                                      & 3. El equipo de Datos (\textit{Data team}) necesita una incorporación adicional para participar                                                                                                                                                                                                      \\
    \midrule

    \multicolumn{2}{l}{\textbf{Decisión: Iterar}}                                                                                                                                                                                                                                                                                            \\* \nopagebreak
                                      & Reducir a 1 PR cruzado cada 2 semanas (más sostenible)                                                                                                                                                                                                                                               \\
                                      & Sesión de 1h para el \textit{Data team} sobre arquitectura general                                                                                                                                                                                                                                   \\
                                      & Reevaluar en el \textit{Sprint} 20                                                                                                                                                                                                                                                                   \\
\end{longtable}
\endgroup